% !PRINTED 155
\weberpage{155}
Ebenso leicht erhalten wir:

\par
5.\enspace\gesperrt{Durch Zusammensetzung mit der identischen Substitution $J$ wird keine Substitution geändert.}

Die Substitution $J$ wird also bei der Composition als \gesperrt{Einheit} betrachtet, und kann, wo sie mit anderen Substitutionen componirt auftritt, weggelassen werden.

Nun gilt weiter der Satz:

\par
6.\enspace\gesperrt{Zu jeder Substitution $A$ giebt es eine und nur eine inverse Substitution $A^{-1}$, die der Bedingung}

\weq{15}{AA^{-1} = A^{-1}A = J}

\gesperrt{genügt.}

Dieser letzte Satz ergiebt sich aus den Grundformeln der Determinantentheorie, wenn man die Elemente $\alpha_h^{(k)}$ von $A^{-1}$ aus den linearen Gleichungen

\weq{16}{\begin{aligned} \sum^i a_i^{(h)}\alpha_k^{(i)} &{}= 0, & h \gtrless k\\ &{}= 1, & h = k \end{aligned}}

bestimmt, aus denen sich, wenn wir mit $A_h^{(k)}$ die Unterdeterminanten von $|A|$ bezeichnen,

\weq{17}{|A|\,\alpha_h^{(k)} = A_k^{(h)}}

ergiebt, und die das andere System

\weq{18}{\begin{aligned} \sum^i \alpha_i^{(h)} a_k^{(i)} &{}= 0, & h \gtrless k\\ &{}= 1, & h = k \end{aligned}}

zur Folge haben. Die inverse Substitution zu $A$ ist nichts Anderes, als die Auflösung des Gleichungssystems (1) der directen Substitution $A$, so dass aus $(y) = A(x)$ folgt:

\weq{19}{(x) = A^{-1}(y).}

Für die Composition der inversen Substitutionen ergiebt sich aus $ABB^{-1}A^{-1} = J$ der Satz:

\weq{20}{(AB)^{-1} = B^{-1}A^{-1}.}

Sind $A$, $B$, $C$ drei Substitutionen, so ergiebt sich durch Zusammensetzung mit $A^{-1}$ aus jeder der beiden Gleichungen

\[
AB = AC, \quad BA = CA,
\]

dass $B = C$ sein muss. Demnach sind für die Zusammensetzung der Substitutionen die charakteristischen Merkmale für eine Gruppe erfüllt, und es ist also der Inbegriff aller Substitutionen
%CONT>>
